%=======================================================================
%  Standalone extract: Literature Review + Research Methodology
%
%  Built by generate_extract.py, which preprocesses the SAME chapter
%  sources the full dissertation uses (report/sections/*.tex) so the two
%  documents cannot drift apart. Do not edit the chapters here.
%=======================================================================
\documentclass[11pt,a4paper]{article}

\usepackage[T1]{fontenc}
\usepackage[utf8]{inputenc}
\usepackage{mathptmx}
\usepackage[scaled=0.88]{helvet}
\usepackage{amsmath,amssymb}
\usepackage[margin=2.6cm,top=2.8cm,bottom=2.9cm]{geometry}
\usepackage{graphicx}
\usepackage{booktabs}
\usepackage{tabularx}
\usepackage{longtable}
\usepackage{multirow}
\usepackage{float}
\usepackage{xcolor}
\usepackage{enumitem}
\usepackage{microtype}
\usepackage{setspace}
\usepackage{fancyhdr}
\usepackage{titlesec}
\usepackage[font=small,labelfont=bf,labelsep=period,justification=justified]{caption}
\usepackage{tcolorbox}
\usepackage{url}
\usepackage[hidelinks]{hyperref}

\definecolor{accent}{HTML}{0072B2}
\definecolor{warm}{HTML}{D55E00}
\definecolor{inkmuted}{HTML}{4D4D4D}
\definecolor{rule}{HTML}{BFBFBF}
\definecolor{panel}{HTML}{F4F7FA}

\hypersetup{colorlinks=true, linkcolor=accent, citecolor=accent,
            urlcolor=accent, linktoc=all,
            pdftitle={Blockchain-Based Electronic Voting:
                      Literature Review and Research Methodology},
            pdfsubject={MSc Computing Science dissertation --- extract}}

\titleformat{\section}{\normalfont\sffamily\large\bfseries\color{black}}
            {\thesection}{0.8em}{}[\vspace{2pt}\color{rule}\hrule height 0.6pt]
\titleformat{\subsection}{\normalfont\sffamily\normalsize\bfseries}{\thesubsection}{0.7em}{}
\titleformat{\subsubsection}{\normalfont\sffamily\small\bfseries}{\thesubsubsection}{0.6em}{}
\titlespacing*{\section}{0pt}{20pt}{10pt}
\titlespacing*{\subsection}{0pt}{14pt}{6pt}

\pagestyle{fancy}
\fancyhf{}
\renewcommand{\headrulewidth}{0.4pt}
\fancyhead[L]{\footnotesize\sffamily\color{inkmuted}Literature Review and Research Methodology}
\fancyhead[R]{\footnotesize\sffamily\color{inkmuted}\thepage}

\setlength{\parskip}{0.42em}
\setlength{\parindent}{0pt}
\onehalfspacing

\newtcolorbox{finding}[1][]{colback=panel, colframe=accent, boxrule=0pt,
  leftrule=2.6pt, arc=1pt, left=9pt, right=9pt, top=7pt, bottom=7pt,
  fontupper=\normalsize, #1}
\newcommand{\finhead}[1]{\textsf{\bfseries #1}\par\vspace{2pt}}

\newcommand{\Prob}{\mathbb{P}}
\newcommand{\Ex}{\mathbb{E}}
\newcommand{\Bin}{\mathrm{Binomial}}
\newcommand{\Uv}{U_{v}}
\newcommand{\figref}[1]{Figure~\ref{#1}}
\newcommand{\tabref}[1]{Table~\ref{#1}}
\newcommand{\eqnref}[1]{Equation~(\ref{#1})}
\newcommand{\code}[1]{\texttt{\small #1}}

% Generated by generate_report.py -- do not edit by hand.
% Every value here is read from results/ output.

\newcommand{\BallotCeilingPct}{30}
\newcommand{\BallotCostRatio}{15}
\newcommand{\BallotStuffingAbsReduction}{0.325}
\newcommand{\BallotStuffingAcceptedBaseline}{0.498}
\newcommand{\BallotStuffingAcceptedProtected}{0.253}
\newcommand{\BallotStuffingBaseline}{0.442}
\newcommand{\BallotStuffingBaselineMedian}{0.0000}
\newcommand{\BallotStuffingBaselineStd}{0.487}
\newcommand{\BallotStuffingCIHigh}{0.343}
\newcommand{\BallotStuffingCILow}{0.308}
\newcommand{\BallotStuffingChernoffGap}{0.0216}
\newcommand{\BallotStuffingChernoffHolds}{100.0}
\newcommand{\BallotStuffingChernoffValid}{3,604}
\newcommand{\BallotStuffingIntervalWidth}{0.043}
\newcommand{\BallotStuffingMCError}{3.94e-4}
\newcommand{\BallotStuffingMCMaxError}{0.0230}
\newcommand{\BallotStuffingProtected}{0.116}
\newcommand{\BallotStuffingRelReduction}{73.7}
\newcommand{\BestPureMax}{0.423}
\newcommand{\BestPureMeanBaseline}{0.277}
\newcommand{\BestPureMin}{0.316}
\newcommand{\BestPureViableCount}{4}
\newcommand{\BriberyAbsReduction}{0.287}
\newcommand{\BriberyBaseline}{0.350}
\newcommand{\BriberyBaselineMedian}{0.0000}
\newcommand{\BriberyBaselineStd}{0.474}
\newcommand{\BriberyCIHigh}{0.304}
\newcommand{\BriberyCILow}{0.269}
\newcommand{\BriberyChernoffGap}{0.0423}
\newcommand{\BriberyChernoffHolds}{100.0}
\newcommand{\BriberyChernoffValid}{1,088}
\newcommand{\BriberyIntervalWidth}{0.018}
\newcommand{\BriberyMCError}{7.98e-5}
\newcommand{\BriberyMCMaxError}{0.0149}
\newcommand{\BriberyProtected}{0.064}
\newcommand{\BriberyRelReduction}{81.9}
\newcommand{\CertificationEpsilon}{10^{-3}}
\newcommand{\ChainIntegrityHeld}{never violated}
\newcommand{\CoalitionBetterCount}{4}
\newcommand{\CoalitionMeanBaseline}{0.800}
\newcommand{\CoalitionMeanGain}{0.522}
\newcommand{\CoalitionMeanProtected}{0.000}
\newcommand{\CoalitionScenarioCount}{5}
\newcommand{\ConvergenceFinalError}{0.0018}
\newcommand{\ConvergenceMaxSize}{50,000}
\newcommand{\ConvergenceMinSize}{10}
\newcommand{\ConvergenceSettles}{3,000}
\newcommand{\ConvergenceTolerance}{0.01}
\newcommand{\CriticalDetectionBindingCount}{4}
\newcommand{\CriticalDetectionMax}{0.118}
\newcommand{\CriticalDetectionMin}{0.050}
\newcommand{\ElectionTrials}{300}
\newcommand{\ElectorateMax}{4,999}
\newcommand{\ElectorateMin}{200}
\newcommand{\FigureCount}{10}
\newcommand{\FlipBallotStuffingBaseline}{55.0}
\newcommand{\FlipBallotStuffingProtected}{8.3}
\newcommand{\FlipBallotStuffingRelReduction}{84.8}
\newcommand{\FlipBriberyBaseline}{44.0}
\newcommand{\FlipBriberyProtected}{1.3}
\newcommand{\FlipBriberyRelReduction}{97.0}
\newcommand{\FlipSybilBaseline}{54.7}
\newcommand{\FlipSybilProtected}{11.3}
\newcommand{\FlipSybilRelReduction}{79.3}
\newcommand{\GateOpenBaseline}{5}
\newcommand{\GateOpenProtected}{0}
\newcommand{\IdentityCostRatio}{40}
\newcommand{\LowerTailShareMax}{30.6}
\newcommand{\LowerTailShareMin}{20.6}
\newcommand{\MonteCarloTrials}{2,000}
\newcommand{\MostResilient}{sybil}
\newcommand{\ProtectedQ}{0.55}
\newcommand{\ProtectionHeadroom}{4.7}
\newcommand{\RuntimeSeconds}{7.9}
\newcommand{\SFiveBStar}{57,100}
\newcommand{\SFiveBudget}{8,000}
\newcommand{\SFiveExactFlip}{0.000}
\newcommand{\SFiveM}{550}
\newcommand{\SFiveMReachable}{76}
\newcommand{\SFiveN}{5,000}
\newcommand{\SFiveQCritical}{0.000}
\newcommand{\SFiveQVoter}{0.167}
\newcommand{\SFiveTStar}{0}
\newcommand{\SFourBStar}{499,400}
\newcommand{\SFourBudget}{5,000,000}
\newcommand{\SFourExactFlip}{0.316}
\newcommand{\SFourM}{4,758}
\newcommand{\SFourMReachable}{4,750}
\newcommand{\SFourN}{5,000}
\newcommand{\SFourQCritical}{0.051}
\newcommand{\SFourQVoter}{0.167}
\newcommand{\SFourTStar}{5,000}
\newcommand{\SOneBStar}{2,880}
\newcommand{\SOneBriberyPFull}{0.423}
\newcommand{\SOneBudget}{3,000}
\newcommand{\SOneExactFlip}{0.423}
\newcommand{\SOneM}{35}
\newcommand{\SOneMReachable}{34}
\newcommand{\SOneN}{100}
\newcommand{\SOneQCritical}{0.118}
\newcommand{\SOneQVoter}{0.118}
\newcommand{\SOneTStar}{37}
\newcommand{\SThreeBStar}{2,998,650}
\newcommand{\SThreeBallotPFull}{0.085}
\newcommand{\SThreeBudget}{3,000,000}
\newcommand{\SThreeExactFlip}{0.320}
\newcommand{\SThreeM}{19,015}
\newcommand{\SThreeMCertified}{19,515}
\newcommand{\SThreeMReachable}{19,000}
\newcommand{\SThreeN}{100,000}
\newcommand{\SThreeQCritical}{0.050}
\newcommand{\SThreeQVoter}{0.111}
\newcommand{\SThreeTStar}{20,000}
\newcommand{\STwoBStar}{249,480}
\newcommand{\STwoBudget}{250,000}
\newcommand{\STwoExactFlip}{0.328}
\newcommand{\STwoM}{1,943}
\newcommand{\STwoMReachable}{1,937}
\newcommand{\STwoN}{10,000}
\newcommand{\STwoQCritical}{0.072}
\newcommand{\STwoQVoter}{0.118}
\newcommand{\STwoTStar}{2,083}
\newcommand{\ScaleBallotLarge}{0.000}
\newcommand{\ScaleBallotSmall}{0.237}
\newcommand{\ScaleBriberyLarge}{0.000}
\newcommand{\ScaleBriberySmall}{1.000}
\newcommand{\ScaleSybilLarge}{0.702}
\newcommand{\ScaleSybilSmall}{0.623}
\newcommand{\Seed}{42}
\newcommand{\SensitivityAtLeastAsGood}{100}
\newcommand{\SensitivityConfigs}{108}
\newcommand{\SensitivityStrictlyBetter}{100}
\newcommand{\SybilAbsReduction}{0.345}
\newcommand{\SybilAcceptedBaseline}{0.537}
\newcommand{\SybilAcceptedProtected}{0.299}
\newcommand{\SybilBaseline}{0.480}
\newcommand{\SybilBaselineMedian}{0.0593}
\newcommand{\SybilBaselineStd}{0.492}
\newcommand{\SybilCIHigh}{0.362}
\newcommand{\SybilCILow}{0.327}
\newcommand{\SybilCeilingPct}{20}
\newcommand{\SybilChernoffGap}{0.0260}
\newcommand{\SybilChernoffHolds}{100.0}
\newcommand{\SybilChernoffValid}{3,471}
\newcommand{\SybilIntervalWidth}{0.044}
\newcommand{\SybilMCError}{3.95e-4}
\newcommand{\SybilMCMaxError}{0.0305}
\newcommand{\SybilProtected}{0.136}
\newcommand{\SybilRelReduction}{71.8}
\newcommand{\TestCount}{122}
\newcommand{\TotalTrials}{2,500}
\newcommand{\VoterIndifferenceMax}{0.167}
\newcommand{\VoterIndifferenceMin}{0.111}
\newcommand{\WEA}{104}
\newcommand{\WEB}{122,592}
\newcommand{\WEN}{4,458}
\newcommand{\WEP}{1051.88}
\newcommand{\WER}{372.81}
\newcommand{\WES}{981}
\newcommand{\WET}{328}
\newcommand{\WEUvBase}{45.85}
\newcommand{\WEUvProt}{-364.10}
\newcommand{\WEballotPBase}{0.000}
\newcommand{\WEballotPProt}{0.000}
\newcommand{\WEballotmuBase}{59.1}
\newcommand{\WEballotmuProt}{27.2}
\newcommand{\WEchernoffBase}{0.875}
\newcommand{\WEchernoffProt}{not applicable}
\newcommand{\WEdecisionBase}{accept}
\newcommand{\WEdecisionProt}{reject}
\newcommand{\WEexactBase}{0.153}
\newcommand{\WEexactProt}{0.000}
\newcommand{\WEm}{261}
\newcommand{\WEmcBase}{0.160}
\newcommand{\WEmcProt}{0.000}
\newcommand{\WEmuBase}{252.7}
\newcommand{\WEmuProt}{0.0}
\newcommand{\WEpbBase}{0.036}
\newcommand{\WEpbProt}{0.281}
\newcommand{\WEpdBase}{0.274}
\newcommand{\WEpdProt}{0.705}
\newcommand{\WEpsurviveBase}{0.771}
\newcommand{\WEpsurviveProt}{0.483}
\newcommand{\WEqBase}{0.229}
\newcommand{\WEqProt}{0.517}
\newcommand{\WEqbBase}{0.589}
\newcommand{\WEqbProt}{0.364}
\newcommand{\WEqsBase}{0.771}
\newcommand{\WEqsProt}{0.425}
\newcommand{\WEsybilPBase}{1.000}
\newcommand{\WEsybilPProt}{0.000}
\newcommand{\WEsybilmuBase}{549.3}
\newcommand{\WEsybilmuProt}{122.8}
\newcommand{\WEtrial}{79}


\begin{document}

%=======================================================================
\begin{titlepage}
\thispagestyle{empty}
\centering
\vspace*{1.4cm}

{\sffamily\Large University of Glasgow}\\[2pt]
{\sffamily\normalsize School of Computing Science}\\[3.4cm]

{\color{rule}\rule{\textwidth}{0.9pt}}\\[0.7cm]
{\sffamily\LARGE\bfseries Blockchain-Based Electronic Voting}\\[0.42cm]
{\sffamily\Large Literature Review and Research Methodology}\\[0.7cm]
{\color{rule}\rule{\textwidth}{0.9pt}}\\[1.4cm]

{\large Mathematical modelling and experimental evaluation of}\\[3pt]
{\large bribery, Sybil and ballot-stuffing attacks}\\[2.8cm]

{\sffamily\large MSc Computing Science}\\[0.35cm]
{\sffamily\normalsize Dissertation --- extract}\\[2.2cm]

\begin{minipage}{0.82\textwidth}
\centering
{\itshape Author, supervisor and submission date to be completed by the candidate.}
\end{minipage}

\vfill
\begin{minipage}{0.86\textwidth}
{\footnotesize\color{inkmuted}
This document contains two chapters of the dissertation --- the
literature review and the research design and methodology --- together
with a short summary of the formal model they refer to. The results
themselves are reported in the full dissertation. References to
chapters not reproduced here are named rather than numbered.
Experimental figures quoted in the methodology are injected from the
accompanying simulation output (seed \Seed, \TotalTrials{} paired
trials) and are reproducible bit-for-bit. No real election system,
voter record or personal datum is used anywhere in this work.}
\end{minipage}
\end{titlepage}

%=======================================================================
\clearpage
\tableofcontents
\clearpage

%=======================================================================
\section{Preliminaries: the model in brief}
\label{sec:prelim}

This section states just enough of the formal model for the two
chapters that follow to be read on their own. It is a summary of the
dissertation's model chapter, not a substitute for it: the derivations,
edge-case conventions and correctness arguments are given there.

\subsection{Threat model}

A single attacker attempts to change the outcome of a two-candidate
election with $N$ enrolled voters. The \emph{manipulation threshold}
$m$ is the number of additional fraudulent votes needed to overturn the
honest margin,
\begin{equation}
m = (\text{opponent votes} - \text{target votes}) + 1 ,
\label{eq:threshold}
\end{equation}
and the attack succeeds precisely when at least $m$ fraudulent votes
survive every countermeasure and reach the final count. Three vectors
are available: \textbf{bribery} of genuine voters, constrained by a
budget $B$ at a reward $R$ per voter and detected with probability $q$;
\textbf{Sybil registration} of $S$ forged identities, each accepted
with probability $q_s$ and detected with probability $p_d$; and
\textbf{ballot stuffing} of $A$ fraudulent ballots, each accepted with
probability $q_b$ and detected with probability $p_b$.

\subsection{The equations referred to below}

A voter offered a bribe faces a decision under risk, gaining $R$ if
undetected and paying a penalty $P$ if not:
\begin{equation}
\Uv = (1 - q) R - q P ,
\label{eq:m1}
\end{equation}
and a risk-neutral voter accepts when $\Uv > 0$. This acts as a
\emph{gate}: where bribery is not incentive-compatible, no bribe is
accepted and the attack contributes nothing at any budget.

Attempts within a vector are independent Bernoulli trials, so the count
of surviving fraudulent votes is binomial and the success probability
is its survival function:
\begin{equation}
\Prob(W \geq m) = \sum_{k=m}^{n} \binom{n}{k} p^{k}(1-p)^{n-k},
\qquad \Ex[W] = \mu = np .
\label{eq:m5}
\end{equation}

A forged identity or an injected ballot must clear two independent
gates to be counted --- acceptance and non-detection --- so the
surviving count is
\begin{equation}
W \sim \Bin\big(n,\; q_{\text{accept}}(1 - p_{\text{detect}})\big).
\label{eq:survival}
\end{equation}

Finally, the attacker's expected utility from $T$ bribery attempts,
given a payoff $V$ for flipping the outcome and a penalty $D$ if the
attack is noticed at all, is
\begin{equation}
U(T) = \Prob(X \geq m)\, V \;-\; T R \;-\; \big[1 - (1-q)^{T}\big] D .
\label{eq:utility}
\end{equation}

Throughout, $q$ denotes the \emph{detection} probability, so a bribed
vote survives with probability $1-q$.

%=======================================================================
\clearpage
%=======================================================================
%  Literature Review
%=======================================================================
\section{Literature Review}
\label{sec:litreview}

\subsection{Scope and method of the review}

This review covers four bodies of work that the dissertation draws
together, and it is organised by the question each body answers rather
than chronologically.

\begin{enumerate}[leftmargin=1.6em,itemsep=2pt]
  \item \textbf{What has been built, and what is claimed for it?}
        Electronic and blockchain-based voting systems, from
        cryptographic verifiable-voting schemes to ledger-backed
        proposals (Sections~\ref{sec:lit-evoting}--\ref{sec:lit-bcvoting}).
  \item \textbf{What breaks when it is deployed?} The empirical
        security literature: analyses of fielded internet-voting systems
        and the critical response to blockchain voting specifically
        (Section~\ref{sec:lit-critique}).
  \item \textbf{How is manipulation modelled formally?} Bribery as a
        computational problem, coercion and receipt-freeness, Sybil
        identity attacks, and the statistical forensics of fraudulent
        tallies (Sections~\ref{sec:lit-bribery}--\ref{sec:lit-forensics}).
  \item \textbf{What tools exist for reasoning about a rational,
        resource-bounded attacker?} Concentration inequalities and
        game-theoretic security analysis
        (Section~\ref{sec:lit-probgame}).
\end{enumerate}

Sources were located by targeted search of the primary security venues
(USENIX Security, ACM CCS, Financial Cryptography, IEEE S\&P and their
workshops), the artificial-intelligence and economics literature for
the manipulation-theory strand, and two recent systematic reviews of
blockchain voting~\cite{ohize2025,rahul2025} used as entry points into
the applied literature. Priority was given to peer-reviewed primary
sources; where a widely-cited result exists only as a technical report
or research blog post, as with the Dark DAO
construction~\cite{daian2018}, this is stated explicitly rather than
disguised by citation formatting.

Section~\ref{sec:lit-synthesis} closes the review by identifying what
the four bodies of work leave unaddressed, which is the gap this
dissertation occupies.

\subsection{Electronic voting before blockchain}
\label{sec:lit-evoting}

The security requirements for an electronic election were substantially
settled before distributed ledgers entered the conversation, and it is
worth restating them because blockchain proposals are frequently
evaluated against a narrower list.

\textbf{Verifiability.} A voter should be able to confirm that their
ballot was recorded as cast, and any observer should be able to confirm
that the recorded ballots were tallied correctly. Adida's
Helios~\cite{adida2008} is the canonical open-audit implementation: it
publishes encrypted ballots to a public bulletin board and proves the
tally correct with homomorphic encryption and zero-knowledge proofs, so
that correctness rests on mathematics rather than on trusting the
election authority. Helios matters to the present work as a
counterexample to a common framing: it delivers public verifiability
with no ledger, no chain and no consensus mechanism, which shows that
the verifiability argument for blockchain voting is not an argument for
blockchain specifically.

\textbf{Ballot secrecy and its stronger forms.} Secrecy alone is
insufficient if a voter can \emph{prove} how they voted, because a
provable vote is a sellable vote. Benaloh and
Tuinstra~\cite{benaloh1994} formalised \emph{receipt-freeness}: the
requirement that a voter cannot construct a convincing proof of their
choice even if they want to. Juels, Catalano and
Jakobsson~\cite{juels2005} strengthened this to
\emph{coercion-resistance}, covering an adversary who can interact with
the voter during voting, demand credentials, or require abstention;
their scheme lets a coerced voter surrender a fake credential
indistinguishable from a real one.

This distinction sets the boundary of the present work. Bribery, as
modelled here, presupposes that the briber can verify compliance well
enough for the transaction to be worth making. A system that is
genuinely receipt-free undermines that presupposition at the root,
which is a different and stronger defence than the detection-based one
modelled in this dissertation. The \emph{Limitations} chapter of the full dissertation returns to this.

\subsection{Blockchain-based voting proposals}
\label{sec:lit-bcvoting}

Nakamoto's chained proof-of-work construction~\cite{nakamoto2008}
supplied a mechanism with obvious surface appeal for elections: an
append-only, hash-linked, replicated log in which retrospective edits
invalidate every subsequent hash and are therefore detectable by anyone
holding a copy.

The most rigorous academic instantiation is McCorry, Shahandashti and
Hao's Open Vote Network~\cite{mccorry2017}, a boardroom-voting protocol
implemented as an Ethereum smart contract, which achieves maximum voter
privacy in the sense that a vote can only be revealed by full collusion
of all remaining participants, and which is self-tallying: no trusted
authority computes the result. It is a genuine cryptographic
contribution rather than a ledger with ballots written into it. Its
scope, however, is explicitly the boardroom --- tens of voters, all
known to one another --- and its cost and collusion properties do not
transfer to a public election.

The wider applied literature is large and, by the assessment of its own
systematic reviewers, uneven. Ohize et al.~\cite{ohize2025} survey
architectures and trends across the field and identify scalability,
post-quantum security, voter authentication, regulatory frameworks and
cross-chain interoperability as the persistent open problems. Rahul,
Gulia and Gill~\cite{rahul2025} review 55 journal articles and reach a
similar conclusion, noting that security, transparency and trust remain
unresolved worldwide rather than settled by any existing proposal. Two
observations from these reviews shape the present work. First,
\emph{voter authentication} is repeatedly named as an open problem,
which is precisely the Sybil vector modelled here. Second, the surveyed
proposals are overwhelmingly evaluated on functionality, throughput and
cryptographic properties rather than on quantified resistance to a
budgeted attacker --- the gap this dissertation addresses.

\subsection{The critical response: what happens on deployment}
\label{sec:lit-critique}

Against the proposal literature stands a smaller and considerably more
sobering body of empirical analysis.

Springall et al.~\cite{springall2014} examined Estonia's I-voting
system, the most established national internet-voting deployment, by
observing operations directly and reconstructing the system in the
laboratory. They reported serious procedural lapses and architectural
weaknesses sufficient, in their assessment, for a state-level attacker
to alter results undetectably, and recommended discontinuing the
system. The finding relevant here is not the specific vulnerabilities
but the location of them: the client platform and the operational
procedures surrounding the count, not the cryptography.

Specter, Koppel and Weitzner~\cite{specter2020} reverse-engineered
Voatz, the first internet-voting application used in U.S. federal
elections and marketed on its use of a blockchain. They found flaws
allowing an adversary to alter, stop or expose a voter's ballot, a
side-channel through which a passive network observer could recover the
secret ballot, and a dependence on third-party services with attendant
privacy exposure; West Virginia and a Washington county withdrew the
application from 2020 primaries following the analysis. The
architectural point is the important one: the blockchain sat behind an
application whose compromise made the ledger irrelevant, because votes
were corrupted before they reached it.

Park, Specter, Narula and Rivest~\cite{park2021} generalise these
observations into a direct argument that blockchain voting is worse
than internet voting rather than better. Their case is that blockchains
address a problem elections do not have --- disagreement about which
transactions occurred among mutually distrusting parties --- while
introducing new failure modes, and that the genuinely hard problems of
elections, namely software independence, secret ballots, voter
authentication and client-side integrity, are untouched. They further
note that a blockchain vote is a \emph{transparent} record, which is in
tension with ballot secrecy rather than supportive of it.

This dissertation adopts that critique as a premise to be tested
quantitatively rather than as a conclusion to be restated. The
experimental design deliberately grants the ledger everything its
proponents claim --- integrity is verified continuously and, as the \emph{Findings} chapter reports, was never violated in any trial ---
and then measures how many elections are stolen anyway.

\subsection{Bribery and vote buying as a formal problem}
\label{sec:lit-bribery}

Three distinct traditions model vote buying, and this dissertation
borrows from all three.

\textbf{Computational.} Faliszewski, Hemaspaandra and
Hemaspaandra~\cite{faliszewski2009} formalise bribery as a decision
problem: given per-voter prices and a budget, can the briber make a
preferred candidate win? They classify the complexity across voting
rules, establishing that some systems make optimal bribery
NP-hard --- complexity as a defensive barrier. The framing is
worst-case, combinatorial and deterministic: a bribed voter votes as
paid.

\textbf{Economic.} Becker's analysis of crime as a decision under
risk~\cite{becker1968} models an offender weighing gain against the
probability of detection multiplied by the penalty, and derives the
policy conclusion that deterrence can be achieved either by raising
detection or by raising penalties. Equation~(\ref{eq:m1}) in this
dissertation is exactly that calculation applied to a bribed voter, and
the critical-detection threshold developed in this work is the
policy question in Becker's sense: how much detection is enough?

\textbf{Blockchain-specific.} Daian, Kell, Bentov and
Juels~\cite{daian2018} describe on-chain vote buying and the
\emph{Dark DAO}: a smart contract that pays for votes and enforces the
payment trustlessly, using trusted execution to verify compliance
without revealing the buyer. Their argument deserves emphasis in the
present context because it inverts the usual claim. A blockchain does
not merely fail to prevent vote buying; by making votes cryptographically
attributable and payments programmable and escrowed, it \emph{improves}
the market for votes, removing the counterparty risk that limits bribery
in traditional settings. This is a technical-report and research-blog
contribution rather than a peer-reviewed paper, and is cited as such,
but it is widely referenced and directly motivates treating bribery,
not tampering, as the primary threat to a ledger-based election.

The gap across all three: the computational tradition treats bribery
deterministically, the economic tradition treats the individual decision
without an electorate-level outcome model, and the blockchain-specific
work is a mechanism proposal rather than a quantified risk model. None
produces the quantity a defender needs --- the probability that a given
budget overturns a given margin against a given detection capability.

\subsection{Sybil attacks and identity}
\label{sec:lit-sybil}

Douceur~\cite{douceur2002} established that, absent a trusted
certification authority, a single entity in a distributed system can
present arbitrarily many distinct identities, and that no purely local
resource test reliably distinguishes them. For an election the
consequence is immediate and fatal: one-person-one-vote is the entire
premise, and a Sybil attack attacks the premise rather than the
mechanism.

Responses divide into those that reintroduce a trusted authority and
those that attempt to establish uniqueness without one. Borge et
al.~\cite{borge2017} propose proof-of-personhood, binding digital
identities to physical presence at pseudonym parties, explicitly to
recover one-per-person semantics in permissionless systems. The
approach illustrates both the necessity and the difficulty: physical
attendance is a strong Sybil barrier and an equally strong barrier to
accessibility, which is what electronic voting exists to improve.

This dissertation does not model any specific identity mechanism. It
parameterises the outcome instead --- a forged registration passes with
probability $q_s$ and is caught with probability $p_d$ --- which is
weaker than assuming a trusted authority and is deliberately
conservative. It also produces the finding, reported in the \emph{Findings} chapter, that identity verification rather than
ledger integrity is where the residual risk concentrates, consistent
with voter authentication being repeatedly flagged as an open problem in
the systematic reviews~\cite{ohize2025,rahul2025}.

\subsection{Ballot stuffing and statistical election forensics}
\label{sec:lit-forensics}

Ballot stuffing --- injecting fraudulent ballots not corresponding to
any voter --- is the least cryptographically interesting of the three
vectors and, historically, among the most common. The relevant
literature is statistical rather than cryptographic.

Klimek, Yegorov, Hanel and Thurner~\cite{klimek2012} show that
systematic irregularities leave detectable signatures in the joint
distribution of turnout and vote share across polling units: incremental
stuffing and extreme fraud produce characteristic distortions absent
from clean elections, and the authors identify these fingerprints in
several national elections. The methodological lesson carried into this
dissertation is that fraud is bounded by \emph{plausibility} as much as
by opportunity. An attacker who injects ballots without limit produces a
turnout figure that requires no probabilistic detection to notice. This
is the justification for the capacity ceilings imposed in the coalition
model (described in the methodology below), which are otherwise a
modelling choice without grounding: they stand in for the anomaly
detection that election forensics performs.

\subsection{Probabilistic and game-theoretic security analysis}
\label{sec:lit-probgame}

The final strand supplies the analytical machinery.

\textbf{Concentration inequalities.} Chernoff's
bound~\cite{chernoff1952}, in the multiplicative form given by
Mitzenmacher and Upfal~\cite{mitzenmacher2017}, bounds the probability
that a sum of independent Bernoulli variables deviates from its mean,
and does so exponentially in the deviation. For this dissertation its
value is specific: it converts an empirical statement into a provable
one. A margin certified safe by the bound is safe irrespective of how
many simulation runs were performed, which is what makes the certified
margin defined in this work a security claim rather than an
observation.

\textbf{Rational attackers.} Eyal and Sirer~\cite{eyal2014}
demonstrated with selfish mining that a blockchain protocol assumed
secure under an honest-majority model can be broken by a rational
minority that deviates profitably --- security depending on incentives,
not only on cryptography or on hash-power share. Liu et
al.~\cite{liu2019} survey the resulting body of game-theoretic
blockchain analysis, covering mining strategy, security games and
incentive design.

Two points transfer directly. First, the relevant question about an
attacker is what they would \emph{choose} to do, not what they are
capable of doing: the \emph{Findings} chapter shows that a
budget-exhausting attacker performs strictly worse than a rational one,
so a defence tuned against maximum expenditure is tuned against the
wrong opponent. Second, deterrence is a legitimate security goal
distinct from prevention --- making an attack unprofitable is cheaper
than making it impossible, and the gap between those two thresholds is
quantified in this dissertation.

\subsection{Synthesis and research gap}
\label{sec:lit-synthesis}

\tabref{tab:gap} summarises what each strand supplies and what it
leaves open.

\begin{table}[htbp]
\centering
\caption{What the literature supplies, and what it leaves open.}
\label{tab:gap}
\small
\begin{tabular}{@{}p{3.5cm}p{5.1cm}p{5.1cm}@{}}
\toprule
\textbf{Strand} & \textbf{Supplies} & \textbf{Leaves open} \\
\midrule
Blockchain voting proposals
  \cite{mccorry2017,ohize2025,rahul2025}
& Architectures, cryptographic protocols, throughput evaluation
& No quantified success probability against a budgeted attacker \\
\addlinespace
Empirical critiques
  \cite{springall2014,specter2020,park2021}
& Evidence that deployments fail upstream of the ledger
& Qualitative argument; no model of how much manipulation gets through \\
\addlinespace
Bribery theory
  \cite{faliszewski2009,becker1968,daian2018}
& Complexity classification; the rational-offender calculus; on-chain vote-buying mechanisms
& Deterministic or individual-level; no electorate-scale outcome distribution \\
\addlinespace
Sybil and identity
  \cite{douceur2002,borge2017}
& Impossibility without trusted certification; countermeasure designs
& No probabilistic model of partial identity verification \\
\addlinespace
Election forensics
  \cite{klimek2012}
& Statistical signatures of stuffing; plausibility as a real constraint
& Detection after the fact, not a design threshold beforehand \\
\addlinespace
Probabilistic and game-theoretic tools
  \cite{chernoff1952,mitzenmacher2017,eyal2014,liu2019}
& Concentration bounds; rational-deviation analysis
& Not yet applied to e-voting manipulation as an integrated model \\
\bottomrule
\end{tabular}
\end{table}

Three gaps follow, and they define the contribution of this
dissertation.

\textbf{Gap 1: no quantified attacker model for e-voting
countermeasures.} The proposal literature evaluates systems on
cryptographic properties and performance; the critical literature
establishes qualitatively that deployments fail. Neither produces the
probability that a specific budget overturns a specific margin against
a specific detection capability. The \emph{Threat Model} and \emph{Findings} chapters supply exactly that, over
\TotalTrials{} matched randomised trials per vector.

\textbf{Gap 2: probabilities are reported where thresholds are
needed.} A defender cannot act on a success-probability curve. The
inverse question --- what detection probability, what budget ceiling,
what margin makes this attack non-viable --- is not asked in the
e-voting literature, although Becker's economics of
deterrence~\cite{becker1968} and the game-theoretic blockchain
work~\cite{eyal2014,liu2019} both imply it should be.
The \emph{Critical Thresholds} chapter formulates seven such thresholds and solves them.

\textbf{Gap 3: attack vectors are evaluated in isolation.} Every source
reviewed treats bribery, Sybil attacks and ballot stuffing separately.
No reviewed work models an attacker allocating one budget across
several vectors simultaneously. The \emph{Coalition Attacks} chapter shows this omission is not conservative: two of the three vectors are individually
infeasible in every benchmark scenario, yet their combination succeeds
decisively in four of five.


%=======================================================================
\clearpage
%=======================================================================
%  Research Design and Methodology
%=======================================================================
\section{Research Design and Methodology}
\label{sec:methodology}

\subsection{Research strategy: why simulation}

The research questions ask what \emph{would} happen to an election
under attack. That rules out the two obvious empirical strategies. A
field experiment is unethical and illegal: it would require attempting
to corrupt a real election. Observational study of historical fraud
suffers from severe selection bias, because only detected fraud is
observable and detected fraud is by definition the fraud that failed at
the part this dissertation measures.

The strategy adopted is therefore \textbf{computational simulation with
analytical validation}: construct a formal model of the attack, compute
its consequences exactly where closed forms exist, simulate where they
do not, and validate the two against each other. The claims that result
are conditional --- ``under this model, with these parameters, the
following follows'' --- and the burden of the design is to make the
conditions explicit, the computation correct, and the results
reproducible.

Three commitments follow, and they shape everything below.

\textbf{Exactness before estimation.} Where a quantity has a closed
form it is computed exactly, not estimated. Attack success probability
is the binomial survival function
(\eqnref{eq:m5}), evaluated through the regularised incomplete beta
function; the coalition distribution is an exact convolution
(Section~\ref{sec:method-coalition}); the threshold solvers are exact
binary searches or bisections on monotone functions. Monte Carlo is
used to \emph{validate} these exact values and to reach quantities that
have no closed form, not as the primary method. This inverts the more
common simulation-study design, and it is what allows the Monte Carlo
error itself to be reported as a result
(see the \emph{Findings} chapter).

\textbf{Falsifiability by construction.} Wherever a modelling decision
could flatter the conclusion, the design includes a mechanism that
would expose it. The Chernoff bound is checked against the exact value
in every trial, so a wrongly-implemented bound would show as a
violation rather than as a plausible number. The coalition convolution
is checked against a 200{,}000-draw Monte Carlo simulation in the test
suite. Reproducibility is enforced by a test that fails if two runs at
the same seed differ, and by a second test that fails if two runs at
\emph{different} seeds agree --- so the first test cannot pass
vacuously.

\textbf{Assumptions declared and swept.} Every modelling assumption is
recorded in \code{ASSUMPTIONS.md} and cross-referenced from the module
that relies on it. Where an assumption is both consequential and
ungrounded --- the relative costs of the three vectors, the capacity
ceilings --- it is not merely declared but swept, and the conclusion is
reported as a function of it (see the \emph{Findings} chapter).

\subsection{A two-layer architecture, and why both layers are necessary}
\label{sec:method-twolayer}

The implementation has two layers that answer the same questions by
different means.

The \textbf{concrete layer} is a working blockchain e-voting
application: SHA-256 block hashing, proof-of-work mining, chain
validation, voter registration, role-based election administration,
duplicate-vote prevention and result publication. Attacking it produces
an actual winner change on an actual mined chain.

The \textbf{abstract layer} is the mathematical framework: exact
binomial evaluation, Chernoff bounds, Monte Carlo, the threshold
solvers and the game-theoretic optimisers. It produces probabilities
rather than outcomes.

Neither layer alone would support the conclusions. The abstract layer
supplies statistical power --- \TotalTrials{} paired trials complete in
\RuntimeSeconds{} seconds, whereas mining that many elections is
infeasible --- but every one of its numbers is a probability computed
under a model, and a model can be wrong in ways that are invisible from
inside it. The concrete layer supplies ground truth --- a real tally, a
real winner, a real chain whose integrity is checked after every attack
--- but at a sample size (\ElectionTrials{} elections) too small for
the parameter sweeps.

Running both and comparing them is a validation strategy, not
duplication. If the abstract model were mis-specified --- if, say, the
survival gates of \eqnref{eq:survival} did not correspond to what
actually happens to a fraudulent ballot in the concrete system --- the
two layers would disagree in ordering or magnitude.
The \emph{Findings} chapter reports that they agree in both.

\subsection{Variables and operationalisation}

\tabref{tab:variables} states how each construct is operationalised.
The distinction that matters most is between the three
\emph{independent} variables the experiment manipulates and the
\emph{control} variables that are held identical across the paired
arms.

\begin{table}[htbp]
\centering
\caption{Operationalisation of the model's variables.}
\label{tab:variables}
\small
\begin{tabular}{@{}p{2.5cm}p{2.6cm}p{8.4cm}@{}}
\toprule
\textbf{Role} & \textbf{Variable} & \textbf{Operationalisation} \\
\midrule
Independent & System configuration & Two levels: baseline (weak detection
and validation) and protected (blockchain-style audit, signatures,
identity verification, ballot validation --- modelled as higher
detection and lower acceptance probabilities). \\
\addlinespace
Independent & Attack vector & Three levels: bribery, Sybil
registration, ballot stuffing; plus the coalition of all three in
the coalition analysis. \\
\addlinespace
Independent & Attacker resources & Budget $B$ and reward $R$, entering
through the feasibility constraint $T = \min(N, \lfloor B/R \rfloor)$. \\
\addlinespace
Dependent & Attack success probability & $\Prob(W \geq m)$, where $W$
counts fraudulent votes surviving every countermeasure. Computed
exactly; also estimated by Monte Carlo and bounded by Chernoff. \\
\addlinespace
Dependent & Election flip & Binary: did the winner of a concretely
mined election change? Measured only in the concrete layer. \\
\addlinespace
Dependent & Attacker expected utility & \eqnref{eq:utility}: payoff
weighted by flip probability, less cost, less expected penalty. \\
\addlinespace
Control & Scenario structure & $N, m, B, R, P, S, A$ --- identical
between the two arms of every pair by construction. \\
\addlinespace
Control & Random stream & A single seeded generator; identical seed
gives identical output bit-for-bit. \\
\bottomrule
\end{tabular}
\end{table}

\subsection{What the ledger contributes, and what it does not}
\label{sec:method-attribution}

Because the protected configuration bundles several controls together,
it is worth stating precisely which component does what before any
result is reported. The pipeline is ordered, and the ordering is the
whole point:

\begin{center}
\small
fraud attempt $\rightarrow$ \textbf{acceptance / identity / signature /
ballot validation} $\rightarrow$ \textbf{detection} $\rightarrow$
surviving votes only $\rightarrow$ \textbf{ledger records them}
\end{center}

Every fraudulent vote that is stopped in this model is stopped by the
validation or detection stage. The ledger sits at the end of the
pipeline and records whatever survives; it contributes integrity,
immutability, auditability and tamper evidence, and it contributes
nothing whatever to deciding \emph{whether a vote should have been
cast}. A blockchain cannot tell that a genuine voter was paid, that a
registration was forged, or that a ballot was injected --- those are
determined upstream, by mechanisms that would work equally well behind
a conventional append-only audit log.

This distinction is preserved in the experimental design rather than
merely asserted in the prose. The concrete layer runs the chain's own
validity check after every attack, so the two properties are measured
separately: the ledger's integrity, and the election's correctness. The
\emph{Findings} chapter reports that these come apart completely ---
integrity held in every trial, and elections were flipped anyway.

Nothing in this dissertation should therefore be read as showing that
``blockchain prevents bribery'', or Sybil attacks, or ballot stuffing.
The claim supported by the results is narrower and conditional: under
the assumptions of this model and over the sampled parameter ranges, a
protected architecture combining ledger-backed auditability with
identity verification, signature validation, ballot validation and
raised detection substantially reduces attack feasibility.

The dependent variable deserves one clarification. ``Attack success''
is defined throughout as reaching the manipulation threshold $m$, which
is the number of additional fraudulent votes needed to overturn the
margin (\eqnref{eq:threshold}). It is not defined as ``some fraud
occurred'' or ``fraud was undetected''. This is a deliberately
demanding definition: an attack that places a hundred fraudulent votes
into an election it fails to flip counts, here, as a failure.

\subsection{Parameter selection and its justification}
\label{sec:method-params}

Parameter ranges (\tabref{tab:parameters}) were selected to span
plausible deployment conditions rather than to be estimated from data,
because for most of these quantities no defensible estimate exists. The
probability that a forged voter registration is detected is not
published by any electoral authority in a form that would support a
point estimate. Three consequences follow, and the design accommodates
all three.

\textbf{Ranges, not points.} Every parameter is sampled from a range on
every trial, so no result depends on a single assumed value. The
electorate spans \ElectorateMin{} to \ElectorateMax{} voters and the
threshold $m$ is itself randomised as 2--25\% of $N$, so the reported
means are averages over a distribution of scenarios rather than
properties of one configuration.

\textbf{Direction of bias stated.} Where a choice could flatter either
side it is made against the conclusion being argued. Legitimate ballots
always pass validation, so the model never credits the defence with
catching fraud by accident. Detection is applied independently per
item, which is more favourable to the attacker than correlated
detection, since an audit that catches one item usually catches its
batch.

\textbf{Sensitivity where it matters.} The coalition analysis depends
on four constants with no empirical basis --- two relative costs and
two capacity ceilings. Rather than defend particular values, the
analysis re-runs across \SensitivityConfigs{} combinations of them and
reports the conclusion's stability
(see the \emph{Findings} chapter).

The benchmark scenarios (S1--S5) are calibrated differently and for a
different purpose. For S1--S4 the threshold $m$ is placed close to the
baseline mean $\mu$ rather than at an extreme. This is necessary for
the three-method comparison to be informative: a threshold far below
$\mu$ makes success trivially certain, one far above makes it trivially
impossible, and the Chernoff bound is not even defined when $m \leq
\mu$. S5 is the deliberate exception --- a threshold far beyond what
its budget can reach --- and functions as a negative control: a model
that reported S5 as succeeding would be wrong in an obvious way.

\subsection{Paired randomised design}
\label{sec:method-paired}

The core experiment is a \textbf{matched-pairs design}. Each trial
draws one scenario and evaluates it twice, under both configurations,
with identical $N, m, B, R, P, S, A$ and identical random streams. Only
the acceptance and detection probabilities differ.

The alternative --- drawing independent samples for each arm --- would
confound the countermeasure effect with the difference between the
scenarios happening to be drawn for each arm. Because the per-trial
success probability has a standard deviation of \BriberyBaselineStd{}
for bribery, scenario-to-scenario variance dominates the effect being
measured, and an unpaired design would need far more trials to reach
the same precision. Pairing removes that variance entirely: the
reported reduction is a within-scenario effect, and every pair is its
own control.

\textbf{Inference.} Confidence intervals are computed by bootstrap
resampling of the paired differences (10{,}000 resamples) rather than
by a normal approximation. This is not a stylistic preference. The
per-trial success probabilities are strongly bimodal --- most scenarios
are either comfortably safe or comfortably lost --- which is visible in
the summary statistics: the baseline bribery distribution has a mean of
\BriberyBaseline{} but a median of \BriberyBaselineMedian. A
distribution with that shape violates the assumptions of a
$t$-interval, and the bootstrap makes no distributional assumption.

\textbf{Trial count.} \TotalTrials{} pairs per vector were run, above
the 1{,}000 originally specified, chosen so that the bootstrap
intervals on the mean reduction are narrow relative to the effect: the
bribery reduction interval $[\BriberyCILow, \BriberyCIHigh]$ has a
width under 0.04 against an effect of \BriberyAbsReduction.

\subsection{Triangulation of three computational methods}
\label{sec:method-triangulation}

Every trial computes its success probability three ways --- exact
binomial, Monte Carlo, and Chernoff bound where valid --- on identical
inputs. This is methodological triangulation, and each method plays a
distinct role.

\begin{description}[leftmargin=1.3em,style=nextline,itemsep=2pt]
\item[Exact binomial.] The ground truth. Every reported probability is
  this value; the other two exist to check it and to characterise their
  own usefulness.
\item[Monte Carlo.] An independent numerical path to the same quantity,
  through a different code path (sampling rather than the survival
  function). Agreement is evidence that neither contains an error;
  its error against the exact value quantifies how much simulation is
  needed for a given precision (RQ2).
\item[Chernoff bound.] A provable upper bound. Its role is not accuracy
  but \emph{certification}: a margin the bound rules out is ruled out
  regardless of simulation effort. Its looseness is measured and
  reported rather than treated as a defect.
\end{description}

The three are computed on the same trial, not on separate runs, which
is what makes the comparison a validation rather than a coincidence.
Critically, the bound is applied only where mathematically valid: the
dispatcher determines whether $m$ lies in the upper tail, the lower
tail, or neither, and returns no bound with a machine-readable reason
outside those regions. Averaging the three quantities across
heterogeneous trials --- as an earlier iteration of this work did ---
compares them on different sub-populations, since the bound exists only
where $m > \mu$; the detection-sweep figure of the \emph{Findings} chapter therefore holds every other
parameter fixed and sweeps only $q$.

\subsection{Threshold solver design}
\label{sec:method-thresholds}

The threshold solvers (T1--T7) invert the model, and three design
decisions govern their trustworthiness.

\textbf{Non-existence is a result.} Each solver returns a structured
object carrying either a value or an explicit absence with a
machine-readable reason. No solver substitutes an endpoint, a zero or a
sentinel for a threshold that does not exist, and the figures print
``does not exist'' rather than drawing a line. A threshold silently
defaulted to a domain boundary would be indistinguishable, in a plot,
from a real finding.

\textbf{Exactness where the domain is discrete.} T3, T6 and T7 search
over integers and use binary search, which is exact because each
underlying function is monotone in the search variable: $\Prob(X \geq
m)$ is non-decreasing in $n$, and the Chernoff bound is decreasing in
$m$ above the mean. These are not approximations.

\textbf{A definitional choice, stated.} T4 and T5 required resolving an
ambiguity in what ``the attacker's utility'' means. The headline
optimiser searches $T \geq 0$ and therefore includes declining to
attack, so its optimum is never negative --- and a quantity floored at
zero cannot be bisected for a zero crossing. The thresholds are
consequently defined on $\max_{T \geq 1} U(T)$, the signed payoff of
actually attacking, which goes negative exactly when attacking loses
money. This is a definition rather than a derivation, it changes the
reported numbers, and it is stated wherever those thresholds appear.

\subsection{Coalition search design}
\label{sec:method-coalition}

The coalition analysis searches the budget simplex exhaustively at a
0.05 step (231 allocations for three vectors) rather than using
continuous optimisation. The reason is structural: the objective is
piecewise flat wherever $\lfloor B/c \rfloor$ does not change, so
gradients are zero almost everywhere and undefined at the steps, which
defeats gradient methods. Exhaustive evaluation is affordable because
each evaluation is a convolution rather than a simulation. The reported
optimum is the best grid point and may differ slightly from the true
continuous optimum --- stated as a limitation rather than glossed.

The success probability of a coalition is computed exactly. A sum of
binomials with unequal success probabilities is not binomial and has no
closed form, and the obvious approximations (normal, Poisson) are least
accurate precisely in the tail that matters. Convolving the component
PMFs is exact, and truncating each at $m-1$ before convolving is also
exact rather than approximate, since $\Prob(W \geq m) = 1 - \Prob(W
\leq m-1)$ and no combination summing to at most $m-1$ can involve a
component larger than $m-1$. That truncation is what makes
national-scale margins tractable: it holds the transform length at $m$
rather than at the total attempt count, which is two orders of
magnitude larger.

\subsection{Validation and verification}
\label{sec:method-validation}

Correctness is established at four levels.

\textbf{Unit level (\TestCount{} tests).} Beyond routine coverage, several tests
encode properties the mathematics requires and that a plausible-looking
implementation can silently violate: that the lower-tail Chernoff bound
carries a negative exponent, since a positive one yields a ``bound''
exceeding 1 that is vacuous but looks reasonable in output; that the
bound is never returned outside its region of validity; that $T = 0$
rather than a division error when $B = 0$, $R = 0$ or $B < R$; and that
each threshold solver's answer actually brackets the crossing it claims
to find.

\textbf{Cross-method level.} The exact and Monte Carlo values are
compared on every trial; the Chernoff bound is checked against the
exact value on every trial where it is defined. A violation would
appear as a bound lying below the exact value it is supposed to bound
(\figref{fig:twosided}); none occurs across \BriberyChernoffValid{},
\SybilChernoffValid{} and \BallotStuffingChernoffValid{} valid trials
respectively.

\textbf{Cross-layer level.} The abstract probabilities and the concrete
mined-election flip rates are compared for ordering and magnitude
(see the \emph{Findings} chapter).

\textbf{Reproducibility level.} A test asserts that two runs at the same
seed produce identical output, and a second asserts that different
seeds do not, so the first cannot pass vacuously.

\subsection{Reproducibility protocol}
\label{sec:reproducibility}

All randomness derives from a single seeded
\code{numpy.random.Generator}; no global random state is used anywhere,
and no unseeded standard-library random calls exist in the codebase.
Re-running with seed \Seed{} reproduces every number in this report
bit-for-bit --- verified by regenerating the full pipeline from a clean
results directory and comparing the summary output. The threshold and
coalition analyses are fully analytic and require no seed at all.

The report itself is part of the protocol. Every number in the prose is
injected at build time from \code{results/json/} and
\code{results/csv/} as a \LaTeX{} macro; none is typed by hand. A macro
whose experiment has not been run causes the document to fail
compilation rather than to print a stale value, which makes it
structurally impossible for the text to disagree with the data it
describes.

\subsection{Ethical considerations}

This work models attacks on elections, which warrants explicit
treatment.

No real system, election, voter or personal datum is involved at any
point. All identifiers are simulation-internal indices; the candidates
are \emph{A} and \emph{B} and have no referent. The artefact is a
mathematical model, not a tool: it computes probabilities and contains
no capability that could be directed at a deployed system, and nothing
in it constitutes an exploit, a technique for evading a specific
control, or an implementation of one.

The disclosure balance favours publication. The vectors modelled ---
bribery, Sybil registration, ballot stuffing --- are long-established
in the literature reviewed in Section~\ref{sec:litreview}, and the
contribution is a quantification of their effect and of the
countermeasure strength required to blunt them. That quantification is
of substantially more use to a defender, who can act on a detection
threshold, than to an attacker, who learns only that the attacks they
already knew about work about as well as expected.


\begin{table}[htbp]
\centering
\caption{Randomised parameter ranges. Every trial draws one value per symbol uniformly from these ranges; the protected configuration modifies only the acceptance and detection probabilities, never $N$, $m$, $B$, $R$, $P$, $S$ or $A$.}
\label{tab:parameters}
\small
\resizebox{\textwidth}{!}{%
\begin{tabular}{llr}
\toprule
Symbol & Meaning & Range \\
\midrule
$N$ & Electorate size & 200 -- 5{,}000 \\
$m$ & Manipulation threshold & 2\% -- 25\% of $N$ \\
$B$ & Attacker budget & 1{,}000 -- 200{,}000 \\
$R$ & Bribe reward per voter & 10 -- 500 \\
$P$ & Penalty if a bribed voter is caught & 10 -- 2{,}000 \\
$q$ & Bribery detection probability (baseline) & 0.02 -- 0.25 \\
$q$ & Bribery detection probability (protected) & baseline $+$ 0.15 -- 0.60 \\
$S$ & Attempted forged identities & 10 -- 1{,}000 \\
$q_s$ & Registration acceptance (baseline) & 0.40 -- 0.95 \\
$p_d$ & Sybil detection (baseline) & 0.02 -- 0.30 \\
$A$ & Attempted fraudulent ballots & 10 -- 1{,}000 \\
$q_b$ & Ballot acceptance (baseline) & 0.30 -- 0.90 \\
$p_b$ & Ballot detection (baseline) & 0.02 -- 0.30 \\
$q_s$ & Registration acceptance (protected) & baseline $\times\,(1 -$ 0.30 -- 0.80$)$ \\
$p_d$ & Sybil detection (protected) & baseline $+$ 0.20 -- 0.70 \\
$q_b$ & Ballot acceptance (protected) & baseline $\times\,(1 -$ 0.30 -- 0.80$)$ \\
$p_b$ & Ballot detection (protected) & baseline $+$ 0.20 -- 0.70 \\
\bottomrule
\end{tabular}%
}
\vspace{2pt}
{\footnotesize\textit{Note.} Source: \texttt{config.py::PARAM\_RANGES}. Currency-like quantities are unitless in the model; only ratios such as $B/R$ affect the results.}
\end{table}


%=======================================================================
\clearpage
\addcontentsline{toc}{section}{References}
\begin{thebibliography}{99}

\bibitem{adida2008} B.~Adida. Helios: Web-based Open-Audit Voting. In
\emph{Proceedings of the 17th USENIX Security Symposium}, pp.~335--348, 2008.

\bibitem{becker1968} G.~S.~Becker. Crime and Punishment: An Economic Approach.
\emph{Journal of Political Economy}, 76(2):169--217, 1968.

\bibitem{benaloh1994} J.~Benaloh and D.~Tuinstra. Receipt-Free Secret-Ballot
Elections (Extended Abstract). In \emph{Proceedings of the 26th Annual ACM
Symposium on Theory of Computing (STOC)}, pp.~544--553, 1994.

\bibitem{borge2017} M.~Borge, E.~Kokoris-Kogias, P.~Jovanovic, L.~Gasser,
N.~Gailly and B.~Ford. Proof-of-Personhood: Redemocratizing Permissionless
Cryptocurrencies. In \emph{IEEE European Symposium on Security and Privacy
Workshops (EuroS\&PW)}, pp.~23--26, 2017.

\bibitem{chernoff1952} H.~Chernoff. A Measure of Asymptotic Efficiency for Tests
of a Hypothesis Based on the Sum of Observations. \emph{Annals of Mathematical
Statistics}, 23(4):493--507, 1952.

\bibitem{daian2018} P.~Daian, T.~Kell, I.~Bentov and A.~Juels. On-Chain Vote
Buying and the Rise of Dark DAOs. \emph{Hacking, Distributed} (IC3 technical
report), July 2018. Available at
\url{https://hackingdistributed.com/2018/07/02/on-chain-vote-buying/}.
\emph{Cited as a technical report; not peer reviewed.}

\bibitem{douceur2002} J.~R.~Douceur. The Sybil Attack. In \emph{Proceedings of
the 1st International Workshop on Peer-to-Peer Systems (IPTPS)}, LNCS 2429,
pp.~251--260, 2002.

\bibitem{eyal2014} I.~Eyal and E.~G.~Sirer. Majority Is Not Enough: Bitcoin
Mining Is Vulnerable. In \emph{Financial Cryptography and Data Security (FC)},
LNCS 8437, pp.~436--454, 2014.

\bibitem{faliszewski2009} P.~Faliszewski, E.~Hemaspaandra and
L.~A.~Hemaspaandra. How Hard Is Bribery in Elections? \emph{Journal of
Artificial Intelligence Research}, 35:485--532, 2009.

\bibitem{juels2005} A.~Juels, D.~Catalano and M.~Jakobsson. Coercion-Resistant
Electronic Elections. In \emph{Proceedings of the 2005 ACM Workshop on Privacy
in the Electronic Society (WPES)}, pp.~61--70, 2005.

\bibitem{klimek2012} P.~Klimek, Y.~Yegorov, R.~Hanel and S.~Thurner.
Statistical Detection of Systematic Election Irregularities. \emph{Proceedings
of the National Academy of Sciences}, 109(41):16469--16473, 2012.

\bibitem{liu2019} Z.~Liu, N.~C.~Luong, W.~Wang, D.~Niyato, P.~Wang,
Y.-C.~Liang and D.~I.~Kim. A Survey on Applications of Game Theory in
Blockchain. \emph{arXiv preprint} arXiv:1902.10865, 2019.

\bibitem{mccorry2017} P.~McCorry, S.~F.~Shahandashti and F.~Hao. A Smart
Contract for Boardroom Voting with Maximum Voter Privacy. In \emph{Financial
Cryptography and Data Security (FC)}, LNCS 10322, pp.~357--375, 2017.

\bibitem{mitzenmacher2017} M.~Mitzenmacher and E.~Upfal. \emph{Probability and
Computing: Randomization and Probabilistic Techniques in Algorithms and Data
Analysis}. 2nd ed., Cambridge University Press, 2017.

\bibitem{nakamoto2008} S.~Nakamoto. \emph{Bitcoin: A Peer-to-Peer Electronic
Cash System}. 2008.

\bibitem{ohize2025} H.~O.~Ohize, A.~J.~Onumanyi, B.~U.~Umar, L.~A.~Ajao,
R.~O.~Isah, E.~M.~Dogo, B.~K.~Nuhu, O.~M.~Olaniyi, J.~G.~Ambafi, V.~B.~Sheidu
and M.~M.~Ibrahim. Blockchain for Securing Electronic Voting Systems: A Survey
of Architectures, Trends, Solutions, and Challenges. \emph{Cluster Computing},
28, 2025. DOI 10.1007/s10586-024-04709-8.

\bibitem{park2021} S.~Park, M.~Specter, N.~Narula and R.~L.~Rivest. Going from
Bad to Worse: From Internet Voting to Blockchain Voting. \emph{Journal of
Cybersecurity}, 7(1):tyaa025, 2021.

\bibitem{rahul2025} Rahul, P.~Gulia and N.~S.~Gill. Articulation of Blockchain
Enabled E-Voting Systems: A Systematic Literature Review. \emph{Peer-to-Peer
Networking and Applications}, 18:142, 2025. DOI 10.1007/s12083-025-01956-3.

\bibitem{specter2020} M.~A.~Specter, J.~Koppel and D.~Weitzner. The Ballot is
Busted Before the Blockchain: A Security Analysis of Voatz, the First Internet
Voting Application Used in U.S. Federal Elections. In \emph{Proceedings of the
29th USENIX Security Symposium}, pp.~1535--1553, 2020.

\bibitem{springall2014} D.~Springall, T.~Finkenauer, Z.~Durumeric, J.~Kitcat,
H.~Hursti, M.~MacAlpine and J.~A.~Halderman. Security Analysis of the Estonian
Internet Voting System. In \emph{Proceedings of the 2014 ACM SIGSAC Conference
on Computer and Communications Security (CCS)}, pp.~703--715, 2014.

\end{thebibliography}

\end{document}
